\begin{lemma}[The Gaussian far majorant for $K_{\gamma,j}$ off the diagonal]
Fix $n \geq 1$ and $j \in \{1,\ldots,n\}$, let $\lambda$ be Lebesgue measure, and let
$K_{\gamma,j}$ be the normalized Gaussian Riesz kernel of \noderef{GaussRieszKernel}, so
that $K_{\gamma,j}(x,y) = c_n\int_0^1\Phi_{x,y}(r)\,dr$ with
$c_n = \sqrt2\,\pi^{-(n+1)/2}$ and
\[
  \Phi_{x,y}(r)
    = \frac{y_j - r x_j}{(1-r^2)^{n/2+1}\sqrt{-\log r}}
      \exp\Bigl(-\frac{|rx-y|^2}{1-r^2}\Bigr),
  \qquad r \in (0,1) .
\]
Then for all $R > 0$ and $\rho > 0$ there is a constant $C > 0$, depending only on $n$,
$j$, $R$ and $\rho$, such that for all $x,y \in \mathbb{R}^n$ with $|x| \leq R$ and
$|x-y| \geq \rho$:
\begin{enumerate}
\item $\Phi_{x,y}$ is Lebesgue integrable on $(0,1)$; and
\item $\bigl|K_{\gamma,j}(x,y)\bigr| \leq C\,e^{-|y|^2}$.
\end{enumerate}

\emph{Why the integrability assertion is part of the content.}
$K_{\gamma,j}(x,y)$ is defined in \noderef{GaussRieszKernel} as a Bochner integral over
$(0,1)$, with the convention that the value is $0$ when the integrand is not integrable;
and $0$ satisfies every upper bound.  A bare size bound
$|K_{\gamma,j}| \leq Ce^{-|y|^2}$ would therefore be provable without establishing
anything about the kernel, and would be useless to a consumer that needs the defining
integral to converge.  Clause (1) is what rules that out, so it is asserted here rather
than assumed.

\emph{The majorant carries no polynomial factor in $|y|$.}
The right-hand side of (2) is $Ce^{-|y|^2}$ with $C$ independent of $y$.  That exact
pairing with the Gaussian density $\pi^{-n/2}e^{-|y|^2}$ is what later lets a truncated
kernel integral be dominated by $\|f\|_{L^1(\gamma_n)}$; a residual factor $|y|^k$ would
break it.  The constant does depend on $R$ and on $\rho$, and no uniformity in $n$, in
$R$ or in $\rho$ is claimed.

The reference paper derives $K_{\gamma,j}$ and fixes the prefactor $c_n$ --- it is
explicit that an unnormalized kernel is a different family of operators, so $c_n$ is not
cosmetic --- but it proves no off-diagonal majorant for it.  Clauses (1) and (2) are
therefore tablet-side rigorization with no paper proof to transcribe: what the paper
supplies is the kernel, and what is needed downstream, for the $L^1$ identification, is
that the kernel is an honest function off the diagonal with Gaussian decay in $y$.
\end{lemma}
\begin{proof}
Fix $n \geq 1$ and $j$, and let $R,\rho > 0$.  Put $R' = \max(R,1) \geq 1$ and
$\rho' = \min(\rho,1) \in (0,1]$; a pair $(x,y)$ with $|x| \leq R$, $|x-y| \geq \rho$ also
satisfies $|x| \leq R'$, $|x-y| \geq \rho'$, so it suffices to produce a constant for
$(R',\rho')$.  Write
\[
  \Psi_{x,y}(r)
    = \frac{y_j - r x_j}{(1-r^2)^{n/2+1}\sqrt{-\log r}}
      \exp\Bigl(-\frac{|x-ry|^2}{1-r^2}\Bigr) ,
  \qquad r \in (0,1) ,
\]
for the reflected integrand.  Throughout, $r \in (0,1)$ gives $1-r^2 > 0$ (because
$r^2 < 1$) and $-\log r > 0$, so every denominator below is nonzero and the prefactor
$\bigl((1-r^2)^{n/2+1}\sqrt{-\log r}\bigr)^{-1}$ is positive.

\emph{Step 0: the splitting identity.}
For all $x,y$ and all $r$,
\begin{equation*}
  |rx-y|^2 - (1-r^2)\bigl(|y|^2-|x|^2\bigr) = |x-ry|^2 , \tag{1}
\end{equation*}
since expanding both sides in the inner product gives
$|x|^2 - 2r\langle x,y\rangle + r^2|y|^2$ for each.  Consequently, for $r \in (0,1)$,
dividing $(1)$ by $1-r^2 > 0$ and rearranging,
\[
  -\frac{|rx-y|^2}{1-r^2}
    = \bigl(|x|^2 - |y|^2\bigr) + \Bigl(-\frac{|x-ry|^2}{1-r^2}\Bigr) ,
\]
and therefore, since $e^{a+b} = e^ae^b$,
\begin{equation*}
  \Phi_{x,y}(r) = e^{|x|^2}e^{-|y|^2}\,\Psi_{x,y}(r)
  \qquad (r \in (0,1)) . \tag{2}
\end{equation*}

\emph{Step 1: the prefactor on $(0,\tfrac12]$.}
For $0 < r \leq \tfrac12$ we have $1-r^2 \geq \tfrac34$ and $r \leq \tfrac12$, so
$\log r \leq \log\tfrac12 = -\log 2$, i.e. $-\log r \geq \log 2 > 0$.  Hence, by
monotonicity of $t \mapsto t^{n/2+1}$ on $[0,\infty)$ and of $\sqrt\cdot$,
\[
  (1-r^2)^{n/2+1}\sqrt{-\log r} \;\geq\; (\tfrac34)^{n/2+1}\sqrt{\log 2} \;>\; 0 ,
\]
so with $K = \bigl((\tfrac34)^{n/2+1}\sqrt{\log 2}\bigr)^{-1} > 0$, a constant depending
only on $n$,
\begin{equation*}
  \frac{1}{(1-r^2)^{n/2+1}\sqrt{-\log r}} \leq K
  \qquad (0 < r \leq \tfrac12) . \tag{3}
\end{equation*}

\emph{Step 2: the bound on $(0,\tfrac12]$.}
First, an elementary inequality: for every $t \geq 0$,
\begin{equation*}
  \bigl((t-R')^+\bigr)^2 \;\geq\; \frac{t^2}{4} - R'^2 , \tag{4}
\end{equation*}
where $s^+ = \max(s,0)$.  If $t \leq 2R'$ the right-hand side is $\leq R'^2-R'^2 = 0$ and
the left-hand side is $\geq 0$; if $t > 2R'$ then $t-R' > t/2 > 0$, so the left-hand side
is $> t^2/4 \geq t^2/4 - R'^2$.

Now fix $x,y$ with $|x| \leq R'$ and let $0 < r \leq \tfrac12$.  From
$|ry| - |x| \leq |x-ry|$ and $|x-ry| \geq 0$ we get
$|x-ry| \geq \bigl(r|y| - R'\bigr)^+$, so squaring (both sides nonnegative) and using
$(4)$ at $t = r|y| \geq 0$,
\[
  |x-ry|^2 \;\geq\; \frac{r^2|y|^2}{4} - R'^2 .
\]
Also $0 < 1-r^2 \leq 1$ gives $-|x-ry|^2/(1-r^2) \leq -|x-ry|^2$.  Combining,
\[
  \exp\Bigl(-\frac{|x-ry|^2}{1-r^2}\Bigr) \leq e^{-|x-ry|^2}
    \leq e^{R'^2}\exp\Bigl(-\frac{|y|^2}{4}r^2\Bigr) .
\]
With $(3)$ and the numerator bound $|y_j-rx_j| \leq |y| + R'$ (which holds because
$|v_i| \leq |v|$ for every coordinate, $r \leq 1$ and $|x| \leq R'$), and since all three
factors are nonnegative,
\begin{equation*}
  \bigl|\Psi_{x,y}(r)\bigr|
    \leq K\,(|y|+R')\,e^{R'^2}\exp\Bigl(-\frac{|y|^2}{4}r^2\Bigr)
  \qquad (0 < r \leq \tfrac12) . \tag{5}
\end{equation*}

\emph{Step 3: the bound on $[\tfrac12,1)$.}
By \noderef{GaussRieszKernelEndpointSliceBound} applied to $R'$ and $\rho'$ there is a
constant $C_2 > 0$, depending only on $n$, $j$, $R'$ and $\rho'$, with
\begin{equation*}
  \bigl|\Psi_{x,y}(r)\bigr| \leq C_2
  \qquad\text{whenever } |x| \leq R',\ |x-y| \geq \rho',\ \tfrac12 \leq r < 1 . \tag{6}
\end{equation*}

\emph{Step 4: one majorant on all of $(0,1)$.}
Fix $x,y$ with $|x| \leq R'$ and $|x-y| \geq \rho'$, and set
\[
  B_1 = K\,e^{R'^2}\,5R' + C_2 \geq 0 , \qquad
  B_2 = K\,(|y|+R')\,e^{R'^2} \geq 0 , \qquad
  b = \frac{|y|^2}{4} \geq 0 ,
\]
\[
  M(r) = B_1 + B_2\,e^{-br^2} .
\]
Then
\begin{equation*}
  \bigl|\Psi_{x,y}(r)\bigr| \leq M(r) \qquad (r \in (0,1)) , \tag{7}
\end{equation*}
because for $r \leq \tfrac12$ the bound $(5)$ is exactly $B_2e^{-br^2}$ and $B_1 \geq 0$,
while for $r > \tfrac12$ the bound $(6)$ gives $C_2 \leq B_1$ and $B_2e^{-br^2} \geq 0$.
Since $e^{-br^2} \leq 1$, $(7)$ also gives the cruder constant bound
$|\Psi_{x,y}| \leq B_1+B_2$ on $(0,1)$.

$\Psi_{x,y}$ is Borel measurable on $\mathbb{R}$ (it is built from $r \mapsto r$,
$\log$, $\sqrt\cdot$, the real power $t \mapsto t^{n/2+1}$, the norm, division and
$\exp$, all Borel), and $\lambda\bigl((0,1)\bigr) = 1 < \infty$, so the crude constant
bound makes $\Psi_{x,y}$ Lebesgue integrable on $(0,1)$; by $(2)$, $\Phi_{x,y}$ is a
constant multiple of $\Psi_{x,y}$ and is therefore integrable on $(0,1)$ as well.  This
is clause (1).

\emph{Step 5: the integral of the majorant.}
$M$ is continuous and bounded by $B_1+B_2$ on $(0,1)$, hence integrable there, and by
additivity and $\lambda((0,1)) = 1$,
\[
  \int_0^1 M \,d\lambda = B_1 + B_2\int_0^1 e^{-br^2}\,dr .
\]
We bound the last term by $K\,e^{R'^2}\bigl(5R' + \tfrac52\sqrt\pi\bigr)$.  By
\noderef{GaussianUnitIntegralLinearFactorBound}, applied with $t = |y| \geq 0$ and
the present $R' \geq 0$, and with $b = |y|^2/4$,
\[
  (|y|+R')\int_0^1e^{-br^2}\,dr \leq 5R' + \tfrac52\sqrt\pi ,
\]
the cited node proving this by the case split on whether $|y| \leq 4R'$.  Since
$B_2 = Ke^{R'^2}(|y|+R')$, multiplying by the nonnegative constant $Ke^{R'^2}$ gives
$B_2\int_0^1e^{-br^2}dr \leq Ke^{R'^2}\bigl(5R'+\tfrac52\sqrt\pi\bigr)$.  Adding the
nonnegative missing summand,
\begin{equation*}
  \int_0^1 M\,d\lambda \;\leq\; C_0
    := K\,e^{R'^2}5R' + C_2 + K\,e^{R'^2}\bigl(5R' + \tfrac52\sqrt\pi\bigr) . \tag{8}
\end{equation*}

\emph{Step 6: clause (2).}
By $(7)$, $(8)$ and $|\int g| \leq \int|g|$,
\[
  \Bigl|\int_0^1 \Psi_{x,y}\,d\lambda\Bigr|
    \leq \int_0^1|\Psi_{x,y}|\,d\lambda
    \leq \int_0^1 M\,d\lambda \leq C_0 .
\]
By $(2)$ and linearity of the integral for the constant $e^{|x|^2}e^{-|y|^2}$,
$\int_0^1\Phi_{x,y} = e^{|x|^2}e^{-|y|^2}\int_0^1\Psi_{x,y}$, so by the definition of
$K_{\gamma,j}$ in \noderef{GaussRieszKernel}, and since $c_n > 0$, $e^{|x|^2} > 0$ and
$e^{-|y|^2} > 0$,
\[
  \bigl|K_{\gamma,j}(x,y)\bigr|
    = c_n\,e^{|x|^2}e^{-|y|^2}\Bigl|\int_0^1\Psi_{x,y}\,d\lambda\Bigr|
    \leq c_n\,e^{R'^2}\,C_0\,e^{-|y|^2} ,
\]
using $|x| \leq R'$ and the monotonicity of $\exp$ to replace $e^{|x|^2}$ by $e^{R'^2}$.
Setting $C = c_n\,e^{R'^2}C_0 > 0$, which depends only on $n$, $j$, $R'$ and $\rho'$ and
hence only on $n$, $j$, $R$ and $\rho$, gives clause (2).
\end{proof}
